\documentclass[11pt]{article}
\usepackage[T1]{fontenc}
\usepackage{amsmath,amssymb,geometry}
\geometry{margin=1in}
\title{The linear polynomial defeats the constant bound\\Disproof of Conjecture 00000000010}
\author{ziangni-sys}
\date{October 8, 2026}
\begin{document}
\maketitle
\paragraph{Scope of the disproof.}
The conjecture concerns a square-free integer polynomial $f$ and asserts,
among other clauses, the explicit upper bound
\[
 C_f\leq2\log D_f,
\]
where $C_f$ is its Bateman--Horn local-density product and $D_f$ is the
discriminant (or its absolute value). Neither language imposes a degree
lower bound, excludes discriminant one, or replaces the discriminant
by a truncated quantity such as $\max\{2,|\operatorname{disc}(f)|\}$.
We disprove this upper bound alone; the remaining asymptotic and
optimality clauses need not be decided. In particular, the disproof
uses no convention about dividing by zero in the separate lower bound.

\paragraph{An admissible polynomial and its discriminant.}
Take the actual polynomial $f(X)=X\in\mathbb Z[X]$. It has degree one,
is primitive, has positive leading coefficient and has no repeated root.
In particular it is square-free. The standard root-product definition
of the discriminant of a split degree-$d$ polynomial with leading
coefficient $a_d$ is
\[
 \operatorname{disc}(f)=a_d^{2d-2}
     \prod_{1\leq i<j\leq d}(r_i-r_j)^2.
\]
Here $d=1$, $a_d=1$, and the single root is $r_1=0$. The product over
root pairs is empty, so it equals one. Consequently
$\operatorname{disc}(X)=1$ and $D_X=1$, for either signed or absolute
discriminant. This value is computed from the definition, not assigned
as an auxiliary certificate.

\paragraph{The genuine local-density product.}
For each prime $p$, the reduction of $X$ in $\mathbb F_p[X]$ has exactly
one root, namely zero. Thus its root count $\rho_X(p)$ is one, and the
single-polynomial Bateman--Horn factor is
\[
 \frac{1-\rho_X(p)/p}{1-1/p}
   =\frac{1-1/p}{1-1/p}=1.
\]
The denominator is positive since $p\geq2$. Every actual prime partial
product
\[
 \prod_{p\leq N}\frac{1-\rho_X(p)/p}{1-1/p}
\]
therefore equals one.
These partial products converge to one, proving $C_X=1$ directly from
the stated local-density product. There is no local obstruction and
no unproved prime-counting assertion in this calculation.

\paragraph{Contradiction.}
The asserted upper bound at this admissible polynomial becomes
\[
 1=C_X\leq2\log D_X=2\log1=0,
\]
which is false. Hence the conjecture, with its stated scope, is false.
Adding a degree or discriminant restriction would change that scope;
the present counterexample does not purport to decide such a modified
claim.

\paragraph{Formal verification.}
The Lean project constructs the polynomial $X$ over $\mathbb Z$, proves
its square-freeness and degree, maps it into $\mathbb C[X]$, verifies
its one-root factorization and computes the standard discriminant
product. It counts the actual zeros of its reduction in \texttt{ZMod p}
for every prime, computes each local factor, forms the prime partial
products, proves their limit, and defines the constant by that limit.
Its final theorem disproves the stated logarithmic upper bound.
\end{document}
